\documentclass[10pt,a4paper]{amsart}
\usepackage{amsmath,amssymb,amsthm,graphicx}
% Retained preamble for the Stacks Project chapter {CH} (shared project preamble).
% The chapter file pinned-original-source.tex only carries \input{{preamble}} on its first line,
% so no preamble text lives inside the chapter itself.  The Stacks Project uses one single shared
% preamble for every chapter: the snapshot stages it three times (sources/stacks-chow/preamble.tex,
% sources/stacks-etale-cohomology/preamble.tex and focused-seeds/AGP/sources/stacks-snapshot-20261011/
% preamble.tex) and those copies are identical in size.  The block below is copied mechanically,
% line for line, from sources/stacks-chow/preamble.tex lines 14-18 and 154-217: the Xy-pic setup and
% the native theorem declarations together with the mathematical macros.  Nothing was added, removed
% or redefined; the wrapper class and the amsmath/amssymb/amsthm/graphicx base packages are supplied
% by the assembler.
% For commutative diagrams we use Xy-pic
% Chapter: stacks-spaces-more-morphisms.tex -- More on Morphisms of Spaces.
\usepackage[all]{xy}

% We use 2cell for 2-commutative diagrams.
\xyoption{2cell}

% Theorem environments.
%
\theoremstyle{plain}
\newtheorem{theorem}[subsection]{Theorem}
\newtheorem{proposition}[subsection]{Proposition}
\newtheorem{lemma}[subsection]{Lemma}

\theoremstyle{definition}
\newtheorem{definition}[subsection]{Definition}
\newtheorem{example}[subsection]{Example}
\newtheorem{exercise}[subsection]{Exercise}
\newtheorem{situation}[subsection]{Situation}

\theoremstyle{remark}
\newtheorem{remark}[subsection]{Remark}
\newtheorem{remarks}[subsection]{Remarks}

\numberwithin{equation}{subsection}

% Macros
%
\def\lim{\mathop{\mathrm{lim}}\nolimits}
\def\colim{\mathop{\mathrm{colim}}\nolimits}
\def\Spec{\mathop{\mathrm{Spec}}}
\def\Hom{\mathop{\mathrm{Hom}}\nolimits}
\def\Ext{\mathop{\mathrm{Ext}}\nolimits}
\def\SheafHom{\mathop{\mathcal{H}\!\mathit{om}}\nolimits}
\def\SheafExt{\mathop{\mathcal{E}\!\mathit{xt}}\nolimits}
\def\Sch{\mathit{Sch}}
\def\Mor{\mathop{\mathrm{Mor}}\nolimits}
\def\Ob{\mathop{\mathrm{Ob}}\nolimits}
\def\Sh{\mathop{\mathit{Sh}}\nolimits}
\def\NL{\mathop{N\!L}\nolimits}
\def\CH{\mathop{\mathrm{CH}}\nolimits}
\def\proetale{{pro\text{-}\acute{e}tale}}
\def\etale{{\acute{e}tale}}
\def\QCoh{\mathit{QCoh}}
\def\Ker{\mathop{\mathrm{Ker}}}
\def\Im{\mathop{\mathrm{Im}}}
\def\Coker{\mathop{\mathrm{Coker}}}
\def\Coim{\mathop{\mathrm{Coim}}}

% Boxtimes
%
\DeclareMathSymbol{\boxtimes}{\mathbin}{AMSa}{"02}

%
% Macros for moduli stacks/spaces
%
\def\QCohstack{\mathcal{QC}\!\mathit{oh}}
\def\Cohstack{\mathcal{C}\!\mathit{oh}}
\def\Spacesstack{\mathcal{S}\!\mathit{paces}}
\def\Quotfunctor{\mathrm{Quot}}
\def\Hilbfunctor{\mathrm{Hilb}}
\def\Curvesstack{\mathcal{C}\!\mathit{urves}}
\def\Polarizedstack{\mathcal{P}\!\mathit{olarized}}
\def\Complexesstack{\mathcal{C}\!\mathit{omplexes}}
% \Pic is the operator that assigns to X its picard group, usage \Pic(X)
% \Picardstack_{X/B} denotes the Picard stack of X over B
% \Picardfunctor_{X/B} denotes the Picard functor of X over B
\def\Pic{\mathop{\mathrm{Pic}}\nolimits}
\def\Picardstack{\mathcal{P}\!\mathit{ic}}
\def\Picardfunctor{\mathrm{Pic}}
\def\Deformationcategory{\mathcal{D}\!\mathit{ef}}
\newenvironment{reference}{}{}
\newenvironment{slogan}{}{}
\newenvironment{history}{}{}
\usepackage{cleveref}
\crefname{originallemma}{original lemma}{original lemmas}
\Crefname{originallemma}{Original lemma}{Original lemmas}
\makeatletter
\newlabel{lemma-affine-locally-rel-perfect}{{0DKQ}{source}{}{sourcepointer-1}{}}
\newlabel{lemma-affine-locally-rel-perfect@cref}{{[originallemma][0][]0DKQ}{[1][0][]source}{}{}{}}
\makeatother
\crefname{originallemma}{original lemma}{original lemmas}
\Crefname{originallemma}{Original lemma}{Original lemmas}
\makeatletter
\newlabel{lemma-relative-descend-homomorphisms}{{0DKV}{source}{}{sourcepointer-2}{}}
\newlabel{lemma-relative-descend-homomorphisms@cref}{{[originallemma][0][]0DKV}{[1][0][]source}{}{}{}}
\makeatother
\crefname{originallemma}{original lemma}{original lemmas}
\Crefname{originallemma}{Original lemma}{Original lemmas}
\makeatletter
\newlabel{perfect-lemma-descend-relatively-perfect}{{0DI8}{source}{}{sourcepointer-3}{}}
\newlabel{perfect-lemma-descend-relatively-perfect@cref}{{[originallemma][0][]0DI8}{[1][0][]source}{}{}{}}
\makeatother
\crefname{originallemma}{original lemma}{original lemmas}
\Crefname{originallemma}{Original lemma}{Original lemmas}
\makeatletter
\newlabel{situation-relative-descent}{{0DKU}{source}{}{sourcepointer-4}{}}
\newlabel{situation-relative-descent@cref}{{[originallemma][0][]0DKU}{[1][0][]source}{}{}{}}
\makeatother
\crefname{originallemma}{original lemma}{original lemmas}
\Crefname{originallemma}{Original lemma}{Original lemmas}
\makeatletter
\newlabel{spaces-perfect-lemma-glue}{{08HB}{source}{}{sourcepointer-5}{}}
\newlabel{spaces-perfect-lemma-glue@cref}{{[originallemma][0][]08HB}{[1][0][]source}{}{}{}}
\makeatother
\crefname{originallemma}{original lemma}{original lemmas}
\Crefname{originallemma}{Original lemma}{Original lemmas}
\makeatletter
\newlabel{spaces-perfect-lemma-induction-principle}{{08GP}{source}{}{sourcepointer-6}{}}
\newlabel{spaces-perfect-lemma-induction-principle@cref}{{[originallemma][0][]08GP}{[1][0][]source}{}{}{}}
\makeatother
\title{Relative perfect complexes on a directed limit of spaces descend to a finite stage}
\author{The Stacks Project Authors}
\date{stacks-project:spaces-more-morphisms.tex:local-snapshot-20261011; GNU FDL 1.2 or later}
\begin{document}
\maketitle
\noindent\textit{Editorial note.} Statement span is lines 12498-12503 of sources/stacks-spaces-more-morphisms/pinned-original-source.tex; proof span is lines 12505-12575, the complete author proof (69 body lines).\par\medskip
\noindent\textit{Editorial note.} Navigation stubs: 6 cited-but-not-included labels are supplied as non-load-bearing pointers with their pinned tags where the snapshot tag index has them.\par\medskip
\noindent\textit{Editorial note.} extra\_packages is empty on purpose: the retained preamble already loads xy with options and a second declaration clashes in the assembler wrapper.\par\medskip
\noindent\textit{Editorial note.} Text is retained verbatim, including the author's own typos and placeholders.\par\medskip
\noindent\textit{Editorial note.} The card's statement is relative to the chapter's Situation for relative descent of perfect complexes.\par\medskip
\noindent\textit{Editorial note.} Cross-chapter same-label warning (NOT an occupancy conflict, Chief ruling 2026-10-13 evening): label 'lemma-descend-relatively-perfect' also exists in stacks-perfect — a different proposition in a different chapter file.\par\medskip
\noindent\textit{Allowed standard prerequisites (editorial).}\begin{itemize}
\item Explicit prerequisite boundary. Full faithfulness of the colimit comparison (same chapter, lemma-relative-descend-homomorphisms), the induction principle for bootstrap on derived categories of spaces (Derived Categories of Spaces, lemma-induction-principle), the affine-local characterization of relative perfectness (lemma-affine-locally-rel-perfect), and the schemes case (Derived Categories of Schemes, lemma-descend-relatively-perfect).
\end{itemize}\medskip
\ifdefined\hypertarget\hypertarget{sourcepointer-1}{}\fi
\noindent\textit{Source reference.} Original lemma 0DKQ: sources/stacks-spaces-more-morphisms/pinned-original-source.tex, source lines 12313--12335. Cited by the retained statement or proof but not part of the collected spans; supplied only as a navigation pointer. The argument does not depend on its content.\par\medskip
\ifdefined\hypertarget\hypertarget{sourcepointer-2}{}\fi
\noindent\textit{Source reference.} Original lemma 0DKV: sources/stacks-spaces-more-morphisms/pinned-original-source.tex, source lines 12435--12449. Cited by the retained statement or proof but not part of the collected spans; supplied only as a navigation pointer. The argument does not depend on its content.\par\medskip
\ifdefined\hypertarget\hypertarget{sourcepointer-3}{}\fi
\noindent\textit{Source reference.} Original lemma 0DI8: sources/stacks-spaces-more-morphisms/pinned-original-source.tex, source lines 12575--12575. Cited by the retained statement or proof but not part of the collected spans; supplied only as a navigation pointer. The argument does not depend on its content.\par\medskip
\ifdefined\hypertarget\hypertarget{sourcepointer-4}{}\fi
\noindent\textit{Source reference.} Original lemma 0DKU: sources/stacks-spaces-more-morphisms/pinned-original-source.tex, source lines 12421--12432. Cited by the retained statement or proof but not part of the collected spans; supplied only as a navigation pointer. The argument does not depend on its content.\par\medskip
\ifdefined\hypertarget\hypertarget{sourcepointer-5}{}\fi
\noindent\textit{Source reference.} Original lemma 08HB: sources/stacks-spaces-more-morphisms/pinned-original-source.tex, source lines 12575--12575. Cited by the retained statement or proof but not part of the collected spans; supplied only as a navigation pointer. The argument does not depend on its content.\par\medskip
\ifdefined\hypertarget\hypertarget{sourcepointer-6}{}\fi
\noindent\textit{Source reference.} Original lemma 08GP: sources/stacks-spaces-more-morphisms/pinned-original-source.tex, source lines 12575--12575. Cited by the retained statement or proof but not part of the collected spans; supplied only as a navigation pointer. The argument does not depend on its content.\par\medskip
\begin{lemma}
\label{lemma-descend-relatively-perfect}
In Situation \ref{situation-relative-descent} the category of
$Y$-perfect objects of $D(\mathcal{O}_X)$ is the colimit of the categories
of $Y_i$-perfect objects of $D(\mathcal{O}_{X_i})$.
\end{lemma}

\begin{proof}
For every quasi-compact and quasi-separated object $U_0$ of
$(X_0)_{spaces, \etale}$ consider the condition $P$
that the functor
$$
\colim_{i \geq 0} D_{Y_i\text{-perfect}}(\mathcal{O}_{U_i})
\longrightarrow
D_{Y\text{-perfect}}(\mathcal{O}_U)
$$
is an equivalence where $U = X \times_{X_0} U_0$ and
$U_i = X_i \times_{X_0} U_0$.
We observe that we already know this functor is fully faithful
by Lemma \ref{lemma-relative-descend-homomorphisms}. Thus it suffices to prove
essential surjectivity.

\medskip\noindent
Suppose that $(U_0 \subset W_0, V_0 \to W_0)$ is an elementary
distinguished square in $(X_0)_{spaces, \etale}$
and $P$ holds for $U_0, V_0, U_0 \times_{W_0} V_0$.
We claim that $P$ holds for $W_0$. We will use the notation
$U_i = X_i \times_{X_0} U_0$, $U = X \times_{X_0} U_0$,
and similarly for $V_0$ and $W_0$. We will abusively use the symbol
$f_i$ for all the morphisms $U \to U_i$, $V \to V_i$,
$U \times_W V \to U_i \times_{W_i} V_i$, and $W \to W_i$.
Suppose $E$ is an $Y$-perfect object of $D(\mathcal{O}_W)$.
Goal: show $E$ is in the essential image of the functor.
By assumption,
we can find $i \geq 0$, an $Y_i$-perfect object $E_{U, i}$ on $U_i$,
an $Y_i$-perfect object $E_{V, i}$ on $V_i$, and
isomorphisms $Lf_i^*E_{U, i} \to E|_U$ and $Lf_i^*E_{V, i} \to E|_V$.
Let
$$
a : E_{U, i} \to (Rf_{i, *}E)|_{U_i}
\quad\text{and}\quad
b : E_{V, i} \to (Rf_{i, *}E)|_{V_i}
$$
the maps adjoint to the isomorphisms $Lf_i^*E_{U, i} \to E|_U$
and $Lf_i^*E_{V, i} \to E|_V$.
By fully faithfulness, after increasing $i$,
we can find an isomorphism
$c : E_{U, i}|_{U_i \times_{W_i} V_i} \to E_{V, i}|_{U_i \times_{W_i} V_i}$
which pulls back to the identifications 
$$
Lf_i^*E_{U, i}|_{U \times_W V} \to E|_{U \times_W V} \to
Lf_i^*E_{V, i}|_{U \times_W V}.
$$
Apply Derived Categories of Spaces, Lemma
\ref{spaces-perfect-lemma-glue}
to get an object $E_i$ on $W_i$ and a map $d : E_i \to Rf_{i, *}E$
which restricts to the maps $a$ and $b$ over $U_i$ and $V_i$.
Then it is clear that $E_i$ is $Y_i$-perfect (because being
relatively perfect is an \'etale local property) and that
$d$ is adjoint to an isomorphism $Lf_i^*E_i \to E$.

\medskip\noindent
By exactly the same argument as used in
the proof of Lemma \ref{lemma-relative-descend-homomorphisms}
using the induction principle
(Derived Categories of Spaces, Lemma
\ref{spaces-perfect-lemma-induction-principle})
we reduce to the case where both $X_0$ and $Y_0$
are affine: first work with quasi-compact and quasi-separated objects
in $(Y_0)_{spaces, \etale}$ to reduce to
$Y_0$ affine, then work with quasi-compact and quasi-separated object
in $(X_0)_{spaces, \etale}$ to reduce to $X_0$ affine.
In the affine case the result follows from the case of schemes which is
Derived Categories of Schemes, Lemma
\ref{perfect-lemma-descend-relatively-perfect}.
The translation into the case for schemes is done by
Lemma \ref{lemma-affine-locally-rel-perfect}.
\end{proof}

\end{document}
